\documentclass[10pt,conference]{IEEEtran}

\usepackage[T1]{fontenc}
\usepackage[utf8]{inputenc}
\usepackage{times}
\usepackage{microtype}
\usepackage{graphicx}
\usepackage{booktabs}
\usepackage{array}
\usepackage{tabularx}
\usepackage{enumitem}
\usepackage{xcolor}
\usepackage{hyperref}
\usepackage{listings}
\usepackage{amsmath}
\usepackage{amssymb}

\hypersetup{
    colorlinks=true,
    linkcolor=blue,
    citecolor=blue,
    urlcolor=blue
}

\setlist[itemize]{leftmargin=*, itemsep=2pt, topsep=2pt}
\setlist[enumerate]{leftmargin=*, itemsep=2pt, topsep=2pt}

\lstset{
    basicstyle=\footnotesize\ttfamily,
    breaklines=true,
    frame=single,
    xleftmargin=4pt,
    xrightmargin=4pt
}

\title{
Assignment 2 Report:\\
Domain Adaptation with Small Language Models\\
for Shakespearean Text Understanding
}

\author{
\IEEEauthorblockN{
Group Number: 
}
\IEEEauthorblockA{
CSCI433/933 Machine Learning Algorithms and Applications\\
}
}

\begin{document}
\maketitle

\section{Dataset and Preprocessing}
\label{sec:dataset}

For this project we used the Shakespeare SLM Dataset provided by the
instructor, which includes three plays: \textit{Hamlet},
\textit{Macbeth}, and \textit{Romeo and Juliet}. Each play comes in
two formats --- a full JSON file containing the complete play, and a
scene-level JSONL file where each line is one scene. There are also
utterance-level JSONL files where each line is one speaker turn.

We used both formats. The scene-level files were used to build the
main retrieval index, and the utterance-level files were used to build
a more detailed index for the enhanced RAG system. We did not use any
external material such as summaries from other sources or glossaries.
All metadata used for enrichment (scene summaries and keywords) came
directly from the provided dataset fields.

Table~\ref{tab:dataset-stats} shows a breakdown of how much data we
had after cleaning.

\begin{table}[ht]
\centering
\caption{Dataset statistics after preprocessing.}
\label{tab:dataset-stats}
\begin{tabular}{lccc}
\toprule
Play & Scenes & Scene Chunks & Utterances \\
\midrule
Hamlet           & 20 & 63  & 1,147 \\
Macbeth          & 28 & 88  & 658   \\
Romeo and Juliet & 25 & 78  & 817   \\
\midrule
\textbf{Total} & \textbf{73} & \textbf{229} & \textbf{2,622} \\
\bottomrule
\end{tabular}
\end{table}

The raw utterance files had 3,613 records in total. After removing
991 stage direction records (where \texttt{speaker} was set to
\texttt{STAGE\_DIRECTION}), we were left with 2,622 utterances. Stage
directions were removed because they describe physical actions on stage
rather than spoken dialogue and would not help with text retrieval.

\subsection{Dataset Schema}

Each processed record follows this structure, which matches the format
recommended in the assignment specification:

\begin{lstlisting}
{
  "play":          "Macbeth",
  "act":           2,
  "scene":         2,
  "speaker":       null,
  "text":          "So foul and fair a day I have not seen.",
  "source_id":     "macbeth_2_2",
  "scene_summary": "Macbeth murders Duncan; guilt begins.",
  "keywords":      ["murder", "guilt", "Duncan"]
}
\end{lstlisting}

For scene-level chunks, the \texttt{speaker} field is set to
\texttt{null} because we keep the full scene as one unit rather than
splitting by speaker. For utterance-level chunks, the
\texttt{speaker} field contains the character name (e.g.,
\texttt{"MACBETH"}) and the \texttt{text} field only contains what
that character said. The \texttt{source\_id} field lets us trace any
retrieved chunk back to the exact play, act, and scene it came from.

\subsection{Cleaning and Filtering}

Before building the retrieval index we applied three cleaning steps to
every record:

\begin{enumerate}
    \item \textbf{Stage direction removal.} We used a regular
    expression to remove stage directions written inside square
    brackets, like \texttt{[Exit Ghost.]} or \texttt{[Aside]}.
    These do not represent dialogue and would just add noise to the
    embeddings. For the utterance files we removed entire records
    where speaker was equal to STAGE\_DIRECTION.

    \item \textbf{Whitespace cleanup.} We collapsed any sequences of
    multiple spaces or newlines into a single space and stripped
    whitespace from the start and end of each text field.

    \item \textbf{Short chunk removal.} Scene chunks shorter than 50
    characters after cleaning were dropped. Utterance chunks shorter
    than 10 characters were also dropped. These fragments are too
    short to carry any useful meaning for retrieval.
\end{enumerate}

\subsection{Chunking Strategy}

We used two different chunking strategies depending on which part of
the system we were building.

\textbf{Scene-level chunks} were used for the main RAG index. Each
scene from the dataset becomes one chunk. We chose this over
utterance-level chunks for the main system because a single line of
dialogue on its own rarely makes sense without context. For example,
the line \textit{``To be or not to be''} tells a beginner very little
on its own --- you need to know who said it, in what situation, and
why. A whole scene gives enough context for the system to produce a
useful answer.

The downside is that scene-level chunks are long and may include
content not directly related to a question. To handle this we did two
things. First, any scene longer than 400 words was split into smaller
overlapping chunks with a 50-word overlap so that the model's
512-token limit is not silently hit. Second, we added the
\texttt{scene\_summary} and \texttt{keywords} from the dataset into
the text before embedding, so retrieval works on meaning rather than
just word matching. After all this we ended up with
\textbf{229 scene-level chunks}.

\textbf{Utterance-level chunks} were used for the enhanced RAG index.
Each individual speaker turn becomes its own chunk. This gives much
more precise retrieval and lets the system pinpoint exactly which
character said something in which scene. The trade-off is that
individual lines lack the surrounding context that scene-level chunks
provide. These 2,622 chunks were stored in a ChromaDB vector database,
which also supports filtering by play, act, scene, or speaker.

\textbf{Text enrichment.} For both chunk types we added the
\texttt{scene\_summary} and \texttt{keywords} to the start of each
chunk before embedding it. An example of the enriched format is:

\begin{lstlisting}
[Scene: Guards and Horatio see the ghost and tell Hamlet.]
[Keywords: ghost, Horatio]
BERNARDO. Who's there? FRANCISCO. Nay, answer me...
\end{lstlisting}

This means a question like \textit{``Why does the ghost appear?''} can
match the right scene even if the word \textit{appear} is not in the
dialogue. For utterance chunks we also added the speaker name to the
start so the system can distinguish between different characters.

\section{Responsible Use of Generative AI}
\label{sec:genai}

\textbf{Relevant marks: Responsible and transparent GenAI usage (10 marks).}

Claude (Anthropic, \texttt{claude.ai}) was used as a reference and
consultation tool at key decision points during this project. Rather
than using it to generate complete solutions, we used it to explore
design options, understand trade-offs, and validate our own reasoning
before making implementation decisions. All outputs were independently
verified against our dataset and requirements before being adopted.

\textbf{What tools were used.} Claude (\texttt{claude.ai}) was the
only generative AI tool used by the Data Engineering Lead and
Evaluation Lead throughout the project.

\textbf{What it was used for.} Claude was consulted to explore
trade-offs in chunking strategy, text enrichment approaches, vector
index design, evaluation methodology, and report structure. It was
also used to investigate the distinction between Masked Language
Modeling and Causal Language Modeling when we were deciding on the
correct training approach for the generative component.

\textbf{How we checked the outputs.} All design decisions informed
by Claude were verified empirically against the actual dataset. For
example, the recommendation to use text enrichment was confirmed by
comparing retrieval results with and without prepended metadata on
the instructor-provided questions. Scores from the LLM-as-judge
component were reviewed and adjusted manually before being included
in the report.

\textbf{What we kept, changed, or threw out.} We adopted the
overlapping chunk splitting approach after verifying it prevented
silent truncation at the 512-token limit. We independently decided
to use cosine similarity for the scene-level index and ChromaDB for
the utterance-level index based on our own assessment of corpus scale
and interface requirements. An early embedding script that used
Masked Language Modeling was discarded after we identified it was
unsuitable for a generative task and rewrote it using Causal Language
Modeling. Report sections suggested by Claude were rewritten where
they did not accurately reflect our implementation.

\textbf{How we ensured submitted work reflects our own understanding.}
Every design decision was discussed as a group before implementation.
We did not adopt any approach without being able to explain it
independently. The evaluation questions, scoring rubric, and failure
analysis were developed by us based on our own observations of the
system's behaviour.

The full usage log is in Appendix~\ref{app:genai-log}.

\appendices


\begin{thebibliography}{9}

\bibitem{reimers2019}
N.~Reimers and I.~Gurevych,
``Sentence-BERT: Sentence Embeddings using Siamese BERT-Networks,''
\textit{Proceedings of EMNLP 2019}, 2019.

\bibitem{lewis2020}
P.~Lewis et al.,
``Retrieval-Augmented Generation for Knowledge-Intensive NLP Tasks,''
\textit{Advances in Neural Information Processing Systems}, 2020.

\bibitem{johnson2021}
J.~Johnson, M.~Douze, and H.~J\'{e}gou,
``Billion-Scale Similarity Search with GPUs,''
\textit{IEEE Transactions on Big Data}, vol.~7, no.~3, pp.~535--547,
2021.

\bibitem{radford2019}
A.~Radford et al.,
``Language Models are Unsupervised Multitask Learners,''
OpenAI Blog, 2019.

\end{thebibliography}

\appendices






\label{app:genai-log}
\begin{table*}[p]
\centering
\caption{GenAI usage log.}
\label{tab:genai-log}
\begin{tabularx}{\textwidth}{p{1.5cm} p{3.5cm} p{3.5cm} X}
\toprule
Tool & Prompt or Task & Useful Output & Verification / Modification \\
\midrule
 
Claude &
Explored the trade-offs between scene-level and utterance-level
chunking strategies for RAG retrieval quality &
A structured comparison of both approaches including context
preservation, retrieval noise, and corpus size considerations &
We evaluated both options against the assignment requirements and
selected scene-level chunking as the primary strategy based on
our own assessment of the dataset size and query types \\
 
Claude &
Investigated how text enrichment using metadata fields affects
embedding quality and retrieval accuracy &
Explanation of how prepending scene summaries and keywords to chunk
text improves semantic alignment between queries and retrieved passages &
We verified this empirically by comparing retrieval results with and
without enrichment on the instructor-provided questions and confirmed
improved relevance \\
 
Claude &
Explored options for building a persistent vector index suitable for
offline deployment on a standard laptop &
Overview of FAISS, ChromaDB, and numpy cosine similarity with
trade-offs for each at different corpus scales &
We independently decided to use cosine similarity for the scene-level
index (229 chunks) and ChromaDB for the utterance-level index (2,622
chunks) based on scale and interface requirements \\
 
Claude &
Consulted on best practices for handling the 512-token limit in
transformer-based embedding models during chunking &
Recommendation to use word-level overlap splitting to prevent silent
truncation at chunk boundaries &
We implemented a 400-word maximum chunk size with 50-word overlap and
verified no scene content was being truncated by inspecting chunk
lengths across all three plays \\
 
Claude &
Discussed evaluation design for a multi-system RAG comparison,
including criteria selection and scoring methodology &
Suggested a 1--5 scoring rubric covering correctness, grounding,
retrieval relevance, usefulness, and style quality &
We refined the rubric to better reflect the assignment criteria and
added group-designed questions targeting areas we identified as
potential failure modes in our system \\
 
Claude &
Explored the concept of LLM-as-judge for automated evaluation
scoring in NLP systems &
Explanation of how a language model can be used to generate
first-pass scores for system outputs, with discussion of its
limitations for small models &
We implemented LLM-as-judge as an initial scoring pass and manually
reviewed and adjusted all scores before including them in the report,
treating the automated scores as a starting point only \\
 
Claude &
Consulted on how to structure a reproducible Python evaluation
pipeline that runs multiple systems and exports results &
A design outline for an automated evaluation script covering
system loading, question iteration, result export to CSV and JSON &
We designed and implemented the final evaluation script ourselves
using the outline as a reference, adapting the structure to match
our specific system architecture and output requirements \\
 
Claude &
Discussed report writing conventions for IEEE-format conference
papers, including how to describe preprocessing decisions &
Guidance on appropriate technical depth, citation style, and how
to justify design choices without being overly verbose &
We wrote all sections ourselves and used Claude to check clarity
and structure; all content was verified against our actual
implementation before inclusion \\
 
\bottomrule
\end{tabularx}
\end{table*}
 
\end{document}